% Bagian: turing-machines
% Bab: undecidability
% Subbagian: decision-problem

\documentclass[../../../include/open-logic-section]{subfiles}

\begin{document}

\olfileid[id]{tur}{und}{dec}
\olsection{Masalah Keputusan}

Kita mengatakan bahwa logika orde pertama \emph{dapat diputuskan} jika dan
hanya jika terdapat metode efektif untuk menentukan apakah suatu
!!{sentence} yang diberikan valid. Ternyata tidak ada metode semacam itu:
masalah untuk memutuskan validitas kalimat-kalimat orde pertama tidak dapat
diselesaikan.

Untuk menetapkan hasil negatif yang penting ini, kita membuktikan bahwa
masalah keputusan tidak dapat diselesaikan oleh mesin Turing. Artinya, kita
menunjukkan bahwa tidak ada mesin Turing yang, setiap kali dijalankan dengan
pita yang memuat suatu !!{sentence} orde pertama, pada akhirnya berhenti dan
menghasilkan $1$ atau~$0$ bergantung pada apakah !!{sentence} tersebut valid.
Menurut tesis Church--Turing, setiap fungsi yang dapat dihitung dapat dihitung
oleh mesin Turing. Jadi, jika ``fungsi validitas'' ini dapat dihitung secara
efektif, fungsi itu dapat dihitung oleh mesin Turing. Jika fungsi itu tidak
dapat dihitung oleh mesin Turing, fungsi itu juga tidak dapat dihitung secara
efektif.

Strategi kita untuk membuktikan bahwa masalah keputusan tidak dapat
diselesaikan ialah mereduksi masalah penghentian kepadanya. Maksudnya sebagai
berikut. Kita telah membuktikan bahwa fungsi~$h(e,w)$ yang berhenti dengan
keluaran~$1$ jika mesin Turing yang dideskripsikan oleh~$e$ berhenti pada
masukan~$w$, dan menghasilkan~$0$ jika tidak, tidak dapat dihitung oleh mesin
Turing. Kita akan menunjukkan bahwa jika ada mesin Turing yang memutuskan
validitas kalimat orde pertama, maka ada pula mesin Turing yang menghitung
$h$. Karena $h$ tidak dapat dihitung oleh mesin Turing, tidak mungkin ada
mesin Turing yang memutuskan validitas.

Langkah pertama dalam strategi ini ialah menunjukkan bahwa untuk setiap
masukan~$w$ dan mesin Turing~$M$, kita dapat mendeskripsikan secara efektif
!!a{sentence} $!T(M, w)$ yang merepresentasikan himpunan instruksi~$M$ dan
masukan~$w$, serta !!a{sentence}~$!E(M, w)$ yang menyatakan ``$M$ pada akhirnya
berhenti'', sedemikian sehingga:
\begin{quote}
  $\Entails !T(M, w) \lif !E(M,w)$ jika dan hanya jika $M$ berhenti pada
  masukan~$w$.
\end{quote}
Bagian terbesar bukti kita akan terdiri atas deskripsi kalimat-kalimat
$!T(M, w)$ dan~$!E(M, w)$ tersebut serta verifikasi bahwa
$!T(M, w) \lif !E(M, w)$ valid jika dan hanya jika $M$ berhenti pada
masukan~$w$.

\end{document}
